\begin{table}[H]
\centering
\captionsetup{justification=centering}
\caption{Comparação entre as duas estratégias para a taxa de desocupação - cálculo indireto e estimação direta pelo modelo multivariado}
\label{tab:comptaxa}
{%
\renewcommand{\arraystretch}{0.8}
\scalebox{0.9}{%
\begin{tabular}{lccccc}
\toprule
& \multicolumn{3}{c}{\textbf{Coeficiente de variação médio (\%)}} & \multicolumn{2}{c}{\textbf{\makecell{Diferença relativa média \\ do erro padrão (\%)}}} \\
\cmidrule(lr){2-4} \cmidrule(lr){5-6}
\multicolumn{1}{l}{\textbf{Estrato Geográfico}} & \textbf{\makecell{Estimativa \\ direta}} & \textbf{\makecell{Cálculo \\ indireto}} & \textbf{\makecell{Modelo \\ direto}} & \textbf{\makecell{Cálculo \\ indireto}} & \textbf{\makecell{Modelo \\ direto}} \\
\midrule
01 - Belo Horizonte & 7,69 & 5,68 & 4,90 & 25,04 & 35,86 \\
02 - Entorno e Colar Metropolitano de BH & 7,33 & 5,82 & 4,95 & 19,04 & 31,30 \\
03 - Sul de Minas & 11,16 & 8,35 & 7,00 & 25,26 & 38,19 \\
04 - Triângulo Mineiro & 12,43 & 7,82 & 6,30 & 35,38 & 48,82 \\
05 - Zona da Mata & 10,84 & 10,29 & 8,27 & 4,34 & 23,33 \\
06 - Norte de Minas & 9,16 & 7,30 & 5,33 & 17,44 & 39,83 \\
07 - Vale do Rio Doce & 10,58 & 10,11 & 7,19 & 5,53 & 32,59 \\
08 - Central & 11,16 & 8,39 & 7,37 & 24,07 & 34,91 \\
\midrule
\textbf{Média} & \textbf{10,04} & \textbf{7,97} & \textbf{6,41} & \textbf{19,51} & \textbf{35,60} \\
\bottomrule
\end{tabular}%
}
}
\fonte{Elaboração própria, com base nos dados da PNAD Contínua. Nota: o cálculo indireto obtém a taxa a partir das tendências estimadas do total de desocupados e do total de ocupados, com variância por linearização de Taylor sob a hipótese \(\mathrm{Cov}(\hat{D}_L,T_L)=0\), discutida no texto. As oito primeiras observações foram descartadas do cálculo.}
\end{table}
